% GENERATED FILE. Edit canonical lessons, then run npm run generate:books.
% canonical-source-hash: fnv1a64:7de67b1282d9f957
% canonical-lessons: TA-C534-torkati, TA-C534-camanilai, TA-C534-natuvar, TA-C534-patakkam, TA-C534-vakaiyalar

\chapter{Winning and Losing}
\label{ch:ta-winning-and-losing}
\hlchaptermodality{534}

\begin{chapteropening}
\textbf{By the end of this chapter:} \emph{I can talk about results, referees and medals.}

Complete the last lesson of chapter 534: I can talk about results, referees and medals.
\end{chapteropening}

\section[\texorpdfstring{\ta{தோற்கடி} (tōṟkaṭi)}{தோற்கடி (tōṟkaṭi)}]{\ta{தோற்கடி} (tōṟkaṭi) — to beat a rival; to defeat}

\label{lesson:TA-C534-torkati}



\begin{quote}
Before the new one: say the Tamil for to leap, then the Tamil for to dive.
\end{quote}

\subsection*{You'll want to know: \ta{தோற்கடி}}
\textbf{\ta{தோற்கடி}} — \emph{tōṟkaṭi} — \textquotedblleft{}to beat a rival; to defeat\textquotedblright{}.

\begin{quote}
\ta{எங்கள்} \ta{அணி} \ta{பலமான} \ta{அணியைத்} \ta{தோற்கடித்தது}.
\end{quote}

\begin{quote}
\emph{eṅkaḷ aṇi palamāṉa aṇiyait tōṟkaṭittatu.}
\end{quote}

\begin{quote}
Our team defeated a strong team.
\end{quote}

Past (I …): \ta{தோற்கடித்தேன்}, \emph{tōṟkaṭittēṉ}.

The side beaten takes -\ta{ஐ}. To be the side that loses is \ta{தோற்றுப்போ} (\emph{tōṟṟuppō}).

\begin{cousinweb}
\textbf{\ta{தோற்கடி}} (\emph{tōṟkaṭi}) is usually read as a native compound: \ta{தோல்} (\emph{tōl}), \textquotedblleft{}to be defeated\textquotedblright{} (DEDR 3558), in its stem \ta{தோற்}-, and \ta{அடி} (\emph{aṭi}), \textquotedblleft{}to strike\textquotedblright{}: to strike someone into defeat.
\end{cousinweb}

\begin{grammarlens}[title={What you've built}]
One more word for describing what you lived and felt.
\end{grammarlens}

\subsection*{Your turn}
\begin{itemize}
\raggedright
  \item \emph{Say it:} \emph{tōṟkaṭi}
  \item \emph{Say it:} \emph{tōṟkaṭi}, once more
  \item \emph{Say it:} \emph{eṅkaḷ aṇi palamāṉa aṇiyait tōṟkaṭittatu.} — Our team defeated a strong team.
  \item \emph{Recall:} say the Tamil for to leap, then the Tamil for to dive, then say \emph{tōṟkaṭi} again
\end{itemize}

\begin{tcolorbox}[breakable,colback=teal!4,colframe=teal!35!black,title={Before you move on}]
What does \ta{தோற்கடி} mean? (\textquotedblleft{}To beat a rival; to defeat\textquotedblright{}.) Say it once more.
\end{tcolorbox}
\section[\texorpdfstring{\ta{சமநிலை} (camanilai)}{சமநிலை (camanilai)}]{\ta{சமநிலை} (camanilai) — a draw; balance}

\label{lesson:TA-C534-camanilai}



\begin{quote}
Before the new one: say the Tamil for to dive, then the Tamil for to beat a rival.
\end{quote}

\subsection*{You'll want to know: \ta{சமநிலை}}
\textbf{\ta{சமநிலை}} — \emph{camanilai} — \textquotedblleft{}a draw; balance\textquotedblright{}.

\begin{quote}
\ta{ஆட்டம்} \ta{சமநிலையில்} \ta{முடிந்தது}.
\end{quote}

\begin{quote}
\emph{āṭṭam camanilaiyil muṭintatu.}
\end{quote}

\begin{quote}
The match ended in a draw.
\end{quote}

In a game it is a draw; elsewhere it is balance: \ta{மனச்} \ta{சமநிலை}, balance of mind.

\begin{cousinweb}
\textbf{\ta{சமநிலை}} (\emph{camanilai}) joins \ta{சம} (\emph{cama}), \textquotedblleft{}equal\textquotedblright{}, from Sanskrit \emph{sama} (CDIAL 13173), and the native \ta{நிலை} (\emph{nilai}), \textquotedblleft{}a state\textquotedblright{}, from \ta{நில்}, \textquotedblleft{}to stand\textquotedblright{} (DEDR 3675).
\end{cousinweb}

\begin{grammarlens}[title={What you've built}]
One more word for describing what you lived and felt.
\end{grammarlens}

\subsection*{Your turn}
\begin{itemize}
\raggedright
  \item \emph{Say it:} \emph{camanilai}
  \item \emph{Say it:} \emph{camanilai}, once more
  \item \emph{Say it:} \emph{āṭṭam camanilaiyil muṭintatu.} — The match ended in a draw.
  \item \emph{Recall:} say the Tamil for to dive, then the Tamil for to beat a rival, then say \emph{camanilai} again
\end{itemize}

\begin{tcolorbox}[breakable,colback=teal!4,colframe=teal!35!black,title={Before you move on}]
What does \ta{சமநிலை} mean? (\textquotedblleft{}A draw; balance\textquotedblright{}.) Say it once more.
\end{tcolorbox}
\section[\texorpdfstring{\ta{நடுவர்} (naṭuvar)}{நடுவர் (naṭuvar)}]{\ta{நடுவர்} (naṭuvar) — a referee; an umpire}

\label{lesson:TA-C534-natuvar}



\begin{quote}
Before the new one: say the Tamil for to beat a rival, then the Tamil for a draw.
\end{quote}

\subsection*{You'll want to know: \ta{நடுவர்}}
\textbf{\ta{நடுவர்}} — \emph{naṭuvar} — \textquotedblleft{}a referee; an umpire\textquotedblright{}.

\begin{quote}
\ta{நடுவர்} \ta{விசில்} \ta{ஊதினார்}.
\end{quote}

\begin{quote}
\emph{naṭuvar vicil ūtiṉār.}
\end{quote}

\begin{quote}
The referee blew the whistle.
\end{quote}

From \ta{நடு}, the middle: the one who stands between the two sides.

\begin{cousinweb}
\textbf{\ta{நடுவர்}} (\emph{naṭuvar}) is native Dravidian: \ta{நடு} (\emph{naṭu}), \textquotedblleft{}the middle, fairness\textquotedblright{} (DEDR 3584), with the personal ending -\ta{அர்}.
\end{cousinweb}

\begin{grammarlens}[title={What you've built}]
One more word for describing what you lived and felt.
\end{grammarlens}

\subsection*{Your turn}
\begin{itemize}
\raggedright
  \item \emph{Say it:} \emph{naṭuvar}
  \item \emph{Say it:} \emph{naṭuvar}, once more
  \item \emph{Say it:} \emph{naṭuvar vicil ūtiṉār.} — The referee blew the whistle.
  \item \emph{Recall:} say the Tamil for to beat a rival, then the Tamil for a draw, then say \emph{naṭuvar} again
\end{itemize}

\begin{tcolorbox}[breakable,colback=teal!4,colframe=teal!35!black,title={Before you move on}]
What does \ta{நடுவர்} mean? (\textquotedblleft{}A referee; an umpire\textquotedblright{}.) Say it once more.
\end{tcolorbox}
\section[\texorpdfstring{\ta{பதக்கம்} (patakkam)}{பதக்கம் (patakkam)}]{\ta{பதக்கம்} (patakkam) — a medal}

\label{lesson:TA-C534-patakkam}



\begin{quote}
Before the new one: say the Tamil for a draw, then the Tamil for a referee.
\end{quote}

\subsection*{You'll want to know: \ta{பதக்கம்}}
\textbf{\ta{பதக்கம்}} — \emph{patakkam} — \textquotedblleft{}a medal\textquotedblright{}.

\begin{quote}
\ta{ஓட்டப்} \ta{போட்டியில்} \ta{அவளுக்கு} \ta{வெள்ளிப்} \ta{பதக்கம்} \ta{கிடைத்தது}.
\end{quote}

\begin{quote}
\emph{ōṭṭap pōṭṭiyil avaḷukku veḷḷip patakkam kiṭaittatu.}
\end{quote}

\begin{quote}
She got a silver medal in the running race.
\end{quote}

Gold, silver and bronze: \ta{தங்கப்} \ta{பதக்கம்}, \ta{வெள்ளிப்} \ta{பதக்கம்}, \ta{வெண்கலப்} \ta{பதக்கம்}.

\begin{cousinweb}
\textbf{\ta{பதக்கம்}} (\emph{patakkam}) is borrowed from Sanskrit \emph{padaka}, \textquotedblleft{}a kind of ornament\textquotedblright{} (Monier-Williams), with the Tamil ending -\ta{அம்} and a doubled \ta{க}.
\end{cousinweb}

\begin{grammarlens}[title={What you've built}]
One more word for describing what you lived and felt.
\end{grammarlens}

\subsection*{Your turn}
\begin{itemize}
\raggedright
  \item \emph{Say it:} \emph{patakkam}
  \item \emph{Say it:} \emph{patakkam}, once more
  \item \emph{Say it:} \emph{ōṭṭap pōṭṭiyil avaḷukku veḷḷip patakkam kiṭaittatu.} — She got a silver medal in the running race.
  \item \emph{Recall:} say the Tamil for a draw, then the Tamil for a referee, then say \emph{patakkam} again
\end{itemize}

\begin{tcolorbox}[breakable,colback=teal!4,colframe=teal!35!black,title={Before you move on}]
What does \ta{பதக்கம்} mean? (\textquotedblleft{}A medal\textquotedblright{}.) Say it once more.
\end{tcolorbox}
\section[\texorpdfstring{\ta{வாகையாளர்} (vākaiyāḷar)}{வாகையாளர் (vākaiyāḷar)}]{\ta{வாகையாளர்} (vākaiyāḷar) — a champion}

\label{lesson:TA-C534-vakaiyalar}



\begin{quote}
Before the new one: say the Tamil for a referee, then the Tamil for a medal.
\end{quote}

\subsection*{You'll want to know: \ta{வாகையாளர்}}
\textbf{\ta{வாகையாளர்}} — \emph{vākaiyāḷar} — \textquotedblleft{}a champion\textquotedblright{}.

\begin{quote}
\ta{இந்த} \ta{ஆண்டும்} \ta{அவர்தான்} \ta{சதுரங்க} \ta{வாகையாளர்}.
\end{quote}

\begin{quote}
\emph{inta āṇṭum avartāṉ caturaṅka vākaiyāḷar.}
\end{quote}

\begin{quote}
This year too he is the chess champion.
\end{quote}

\ta{வாகை}, the garland of victory, + -\ta{ஆளர்} \textquotedblleft{}one who …\textquotedblright{}.

\begin{cousinweb}
\textbf{\ta{வாகையாளர்}} (\emph{vākaiyāḷar}) is built on \ta{வாகை} (\emph{vākai}) with -\ta{ஆளர்}, \textquotedblleft{}one who …\textquotedblright{}. DEDR 5333 gives \ta{வாகை} as the name of the sirissa tree; how it came to mean victory is not settled.
\end{cousinweb}

\begin{grammarlens}[title={What you've built}]
One more word for describing what you lived and felt.
\end{grammarlens}

\subsection*{Your turn}
\begin{itemize}
\raggedright
  \item \emph{Say it:} \emph{vākaiyāḷar}
  \item \emph{Say it:} \emph{vākaiyāḷar}, once more
  \item \emph{Say it:} \emph{inta āṇṭum avartāṉ caturaṅka vākaiyāḷar.} — This year too he is the chess champion.
  \item \emph{Recall:} say the Tamil for a referee, then the Tamil for a medal, then say \emph{vākaiyāḷar} again
\end{itemize}

\begin{tcolorbox}[breakable,colback=teal!4,colframe=teal!35!black,title={Before you move on}]
What does \ta{வாகையாளர்} mean? (\textquotedblleft{}A champion\textquotedblright{}.) Say it once more.
\end{tcolorbox}
